\documentclass[12pt]{article}
\usepackage{../../ficha}

\begin{document}

\ficha{Prioridade das Operações}[Propriedades, Cálculo Mental e Expressões Numéricas]

\instrucoes{Resolve sem calculadora, apresentando os passos intermédios. Lembra-te da ordem: parênteses, potências, multiplicações e divisões, adições e subtrações.}

\secao{Propriedades das Operações}{Indica a propriedade usada em cada igualdade (comutativa, associativa, elemento neutro, elemento absorvente ou distributiva).}
\begin{exercicios}
  \item $17 + 25 = 25 + 17$
  \item $(8 + 13) + 7 = 8 + (13 + 7)$
  \item $46 \times 1 = 46$
  \item $0 + 93 = 93$
  \item $125 \times 0 = 0$
  \item $4 \times (9 \times 25) = (4 \times 25) \times 9$
  \item $6 \times (10 + 3) = 6 \times 10 + 6 \times 3$
  \item $12 \times 7 + 12 \times 3 = 12 \times (7 + 3)$
  \item $9 \times (20 - 1) = 180 - 9$
  \item A subtração é comutativa? Dá um exemplo que justifique a tua resposta.
\end{exercicios}

\secao{Cálculo Mental}{Calcula mentalmente, usando as propriedades das operações.}
\begin{exercicios}
  \item $37 + 48 + 63$
  \item $25 \times 17 \times 4$
  \item $5 \times 46 \times 2$
  \item $99 + 256$
  \item $398 + 145$
  \item $15 \times 12$
  \item $49 \times 6$
  \item $101 \times 23$
  \item $125 \times 8 \times 7$
  \item $720 : 8$
  \item $36 \times 25$
  \item $1\,000 - 357$
\end{exercicios}

\secao{Estimativas}{Faz uma estimativa arredondando os números e depois calcula o valor exato.}
\begin{exercicios}
  \item $398 + 503$
  \item $1\,198 - 402$
  \item $49 \times 21$
  \item $31 \times 19$
  \item $8\,100 : 9$
  \item O João comprou 3 livros de 19 euros e 2 jogos de 31 euros. Estima o valor total e depois calcula-o.
  \item Um cinema tem 28 filas de 32 lugares. Estima o número de lugares e depois calcula-o.
  \item A estimativa de $612 : 3$ deve ser próxima de 20, 200 ou 2\,000? Justifica.
\end{exercicios}

\secao{Expressões sem Parênteses}{Calcula o valor de cada expressão.}
\begin{exercicios}
  \item $5 + 3 \times 4$
  \item $20 - 12 : 4$
  \item $18 : 3 + 2 \times 5$
  \item $7 \times 6 - 4 \times 8$
  \item $40 - 6 \times 5 + 2$
  \item $36 : 4 : 3$
  \item $48 : 6 \times 2$
  \item $3 + 2^3 \times 2$
  \item $50 - 4^2 : 2$
  \item $2 \times 3^2 - 5 \times 2 + 1$
  \item $100 : 5^2 + 3 \times 7$
  \item $6^2 - 2^5 + 10 : 2$
\end{exercicios}

\secao{Expressões com Parênteses}{Calcula o valor de cada expressão.}
\begin{exercicios}
  \item $(5 + 3) \times 4$
  \item $(20 - 12) : 4$
  \item $3 \times (8 - 2) + 4$
  \item $(15 + 9) : (10 - 4)$
  \item $2 \times (3 + 4)^2$
  \item $(2 \times 3 + 4)^2$
  \item $60 : (3 \times 4 - 2)$
  \item $[20 - (3 + 5)] \times 2$
  \item $5 \times [12 - 2 \times (7 - 3)]$
  \item $\{40 - [6 + 2 \times (9 - 4)]\} : 4$
  \item $(4^2 - 7) \times (3 + 2^3)$
  \item $[3 \times (2 + 5) - 1] : (2^2 + 1)$
\end{exercicios}

\secao{Desafios}{Resolve os problemas.}
\begin{exercicios}
  \item Coloca parênteses em $6 + 4 \times 5 - 2$ de modo a obter 48.
  \item Coloca parênteses em $24 : 4 + 2 \times 3$ de modo a obter 12.
  \item Coloca parênteses em $3 + 5 \times 8 - 6$ de modo a obter 16.
  \item Escreve uma expressão numérica que traduza: «ao dobro da soma de 7 com 5 subtrai o triplo de 4». Calcula o seu valor.
  \item A Ana comprou 4 cadernos a 3 euros cada e 2 canetas a 2 euros cada. Pagou com uma nota de 20 euros. Escreve uma única expressão numérica para o troco e calcula-o.
  \item Usando quatro vezes o algarismo 4 e quaisquer operações e parênteses, escreve expressões cujo valor seja 0, 1, 2 e 3.
  \item O Rui calculou $8 + 2 \times 5 = 50$. Explica o erro e corrige-o.
  \item Calcula $1 + 2 + 3 + \cdots + 20$ usando a propriedade comutativa e associativa de forma inteligente.
\end{exercicios}

\end{document}
